Apo\documentclass[11pt,a4paper]{article}
\usepackage[T1]{fontenc}
\usepackage[utf8]{inputenc}
\usepackage[lithuanian]{babel}
\usepackage{lmodern}
\usepackage{microtype}
\setlength{\emergencystretch}{1.5em}
\usepackage{amsmath,amssymb,array,tabularx,enumitem,fancyhdr}
\usepackage[a4paper,margin=1.7cm,headheight=15pt]{geometry}
\setlength{\parindent}{0pt}
\setlength{\parskip}{5pt}
\setlist[itemize]{leftmargin=1.5em,itemsep=2pt,topsep=3pt}
\pagestyle{fancy}\fancyhf{}
\lhead{Kvadratinės lygtys: vertinimo instrukcija}\rhead{Mokytojui}
\cfoot{\thepage}
\newcommand{\uzd}[3]{\subsection*{#1 uždavinys \hfill #2 taškai}\textit{Sąlyga.} #3\par}
\newcommand{\kriterijai}{\textbf{Taškų skyrimas:}}
\newcommand{\balas}[2]{\item \textbf{#1 t.} #2}
\begin{document}
\begin{center}
{\LARGE\bfseries Vertinimo instrukcija}\par
{\large Apolono variantas}\par
Varianto maksimumas -- \textbf{40 taškų}.
\end{center}
\textbf{Šaltinis.} Instrukcija parengta pagal paskutinėje žinutėje pateiktas uždavinių sąlygas, atsižvelgiant į vėlesnį Apolono 6 uždavinio patikslinimą. Išlaikyti jose nurodyti uždavinių taškai; jų vidinis paskirstymas sudarytas pagal vertinamus gebėjimus. Lygiaverčiai teisingi sprendimo būdai vertinami tuo pačiu maksimumu.

\section*{Bendrosios taisyklės}
\begin{itemize}
\item Taškai skiriami už parodytą matematinį turinį, ne už formulių ar eilučių kiekį. Sujungti, bet aiškiai pagrįsti žingsniai gauna visus atitinkamus taškus.
\item Vertinama atskirai kiekviename \textbf{savarankiškame vertinimo vienete}: visas uždavinys, o uždaviniuose su dalimis -- kiekviena a), b), c) dalis. Taigi antroji klaida vienoje dalyje nepanaikina teisingai išspręstos kitos dalies taškų. Tai šioje instrukcijoje pasirinkta klaidų taisyklės taikymo apimtis.
\item \textbf{Viena aritmetinė klaida.} Jei po vienos skaičiavimo klaidos mokinys nuosekliai ir matematiškai teisingai tęsia sprendimą su savo gautais skaičiais, atimamas tik \textbf{1 taškas} tame vertinimo vienete. Tos klaidos sukeltas neteisingas galutinis atsakymas papildomai nebaudžiamas.
\item \textbf{Antroji nepriklausoma aritmetinė klaida.} Jei tame pačiame vienete pasitaiko dar viena neištaisyta aritmetinė klaida, skiriami tik taškai už teisingai atliktus kriterijus \textbf{iki pirmosios aritmetinės klaidos}. Už pirmąją klaidą dar vienas taškas papildomai neatimamas. Už vėlesnę sprendimo dalį taškų neskiriama.
\item \textbf{Klaidos pasekmė nėra nauja klaida.} Pavyzdžiui, klaidinga šaknų suma, iš jos teisingai apskaičiuota kvadratų suma ir iš jos teisingai sudaryta lygtis -- viena klaida. Nauja klaida būtų neteisingai suskaičiuota paties mokinio gauta skaičių išraiška.
\item Aiškiai ištaisyta klaida nelaikoma neištaisyta klaida. Vertinamas galutinis, aiškiai pasirinktas sprendimo kelias. Jei pateikti keli prieštaringi atsakymai ir nė vienas nepasirinktas, galutinio atsakymo kriterijus neįvykdytas.
\item \textbf{Metodinė klaida} (neteisinga formulė, neteisėtas dalijimas, neteisingas keitinys, nepagrįstas sprendinių atmetimas) nevertinama kaip vien aritmetinė. Taškai už klaidingą metodą ir nuo jo priklausomas nepagrįstas išvadas neskiriami; nepriklausomi teisingi žingsniai gali būti įskaityti.
\item Vienos klaidos taisyklė nesuteikia taškų už neparodytus gebėjimus: jei klaida uždavinį supaprastino ir mokinys nebetikrino šakų ar sąlygų, neparodyti kriterijai neįskaitomi. Vertinimo minimumas -- 0 taškų.
\end{itemize}

\textbf{Tik atsakymas.} Be sprendimo galima skirti tik aiškiai atsakymui numatytus taškus. Išimtis -- tiesioginis skaidymo dauginamaisiais užrašas: teisinga pilna sandauga pati parodo prašytą rezultatą ir gauna visus tai daliai skirtus taškus.

\newpage
\section*{Kaip praktiškai taikyti klaidų taisyklę}
\textbf{Viena klaida:} nustatomas balas už parodytus teisingus metodus, vertinant vėlesnį skaičiavimą pagal mokinio gautus skaičius; iš šio balo atimamas 1 taškas. To paties klaidingo skaičiaus pasikartojimai ar klaidingas atsakymas antrą kartą nebaudžiami.

\textbf{Dvi klaidos:} grįžtama prie pirmosios klaidos vietos ir susumuojami tik iki jos teisingai įvykdyti kriterijai. Jei kriterijus reikalauja ir metodo, ir rezultato, o rezultatas iki pirmos klaidos dar negautas, už tą kriterijų taškas neskiriamas. Kai eilutėje numatyti 2 ar 3 taškai, jų vidinis paskirstymas žemiau nurodytas atskirai.

\subsection*{Pavyzdys: Apolono 5 uždavinys (5 t.)}
Teisingai reikia gauti $S=x_1+x_2=6$, $P=x_1x_2=4$, naują sumą $S^2-2P=28$ ir naują sandaugą $P^2=16$.
\begin{itemize}
\item Mokinys teisingai užrašo $S=6$, $P=4$, bet apskaičiuoja $36-8=30$. Toliau teisingai gauna naują sandaugą $16$ ir užrašo $y^2-30y+16=0$. \textbf{4/5 t.}: viena aritmetinė klaida, tolesnis metodas teisingas.
\item Tas pats mokinys papildomai užrašo $4^2=12$. \textbf{3/5 t.}: iki pirmos klaidos teisingai nustatė $S$ (1 t.), $P$ (1 t.) ir naujos sumos formulę $S^2-2P$ (1 t.). Taškai už vėlesnį darbą nebeskiriami.
\item Mokinys užrašo $x_1^2+x_2^2=(x_1+x_2)^2$. Tai \textbf{metodinė klaida}: praleistas $-2x_1x_2$ narys. Vieno taško aritmetinės klaidos lengvata čia netaikoma.
\end{itemize}

\subsection*{Kiti svarbūs sprendimo požymiai}
\begin{itemize}
\item Taškas už diskriminantą nėra privalomas sprendžiant skaidymu ar Vijeto teorema: jis pakeičiamas lygiaverčiu teisingu šaknų nustatymo žingsniu.
\item Už teisingas šaknis nereikalaujama dar ir įstatymo į pradinę lygtį, jei visas kelias buvo ekvivalentus ir apibrėžimo sritis nepažeista.
\item Sąlygose „neskaičiuodami sprendinių“ tikslas -- įvertinti darbą su suma ir sandauga. Vien tik aiškiai apskaičiuotų šaknų naudojimas nesuteikia taškų už prašytą šaknų neskaičiavimo metodą; galutinio rezultato taškai gali būti skiriami.
\item Galutinio atsakymo kriterijus suteikiamas vieną kartą. Jei du sprendiniai numatyti viename 1 taško kriterijuje, reikia abiejų; dalinis sprendinių radimas gali būti įskaitytas ankstesniame metodo kriterijuje.
\end{itemize}

\subsection*{Pažymio skyrimas}
\begin{center}\small
\begin{tabular}{|c|*{10}{c|}}\hline
Taškai&0--4&5--8&9--12&13--16&17--20&21--24&25--29&30--33&34--37&38--40\\\hline
Pažymys&1&2&3&4&5&6&7&8&9&10\\\hline
\end{tabular}\end{center}

\textbf{Patikslinta sąlyga.} 6 uždavinio pirmasis daugiklis yra $-x^2+6x-5$. Teisingi lygties sprendiniai: $0$ ir $6$.
\newpage
\section*{Apolono variantas: 1--4 uždaviniai}
\uzd{1}{3}{$(x-2)(x+3)=2x+6$.}
\textbf{Sprendimas.} $x^2+x-6=2x+6\iff x^2-x-12=0\iff(x-4)(x+3)=0$. Atsakymas: $x=-3$ arba $x=4$.
\kriterijai\begin{itemize}
\balas{1}{Teisingai išskleidžia ir gauna $x^2-x-12=0$.}
\balas{1}{Teisingai skaido $(x-4)(x+3)=0$ arba nustato šaknis kitu metodu.}
\balas{1}{Pateikia abu sprendinius $-3$ ir $4$.}
\end{itemize}
\uzd{2}{3}{$x^2+px-10=0$ šaknis $-2$; rasti $p$ ir kitą šaknį.}
\textbf{Sprendimas.} $4-2p-10=0\Rightarrow p=-3$. Iš $(-2)x_2=-10$ gauname $x_2=5$.
\kriterijai\begin{itemize}
\balas{1}{Įstato duotą šaknį: $4-2p-10=0$, arba iš sandaugos pradeda lygiavertį kelią.}
\balas{1}{Nustato $p=-3$.}
\balas{1}{Nustato kitą šaknį $5$.}
\end{itemize}
\textit{Vertinimo pastaba.} Galima pirmiausia rasti $x_2=5$, o tada iš šaknų sumos $p=-3$. Už du vien teisingus atsakymus be sprendimo -- iki 2 t.
\uzd{3}{6}{a) $-3x^2+x+2$; b) $x(x-y-z)+y(z+y-x)$; c) $x^3+2x^2-9x-18$.}
\textbf{Sprendimas.} \textbf{a)} $-3x^2+x+2=-(3x+2)(x-1)$.\par\textbf{b)} $z+y-x=-(x-y-z)$, todėl reiškinys lygus $(x-y)(x-y-z)$.\par\textbf{c)} $x^2(x+2)-9(x+2)=(x+2)(x-3)(x+3)$.
\kriterijai\begin{itemize}
\balas{2}{a): 1 t. už tinkamas šaknis arba tinkamą skaidymo pradžią; 1 t. už pilną sandaugą.}
\balas{2}{b): 1 t. už priešingų skliaustų atpažinimą; 1 t. už $(x-y)(x-y-z)$.}
\balas{2}{c): 1 t. už $(x+2)(x^2-9)$; 1 t. už pilną skaidymą.}
\end{itemize}
\textit{Vertinimo pastaba.} Teisinga pilna sandauga be tarpinių eilučių -- 2 t. už atitinkamą dalį.
\uzd{4}{4}{Stačiakampio plotas $84\ \mathrm{cm}^2$, perimetras $38\ \mathrm{cm}$.}
\textbf{Sprendimas.} Kraštinėms $u,v>0$: $uv=84$, $u+v=19$. Todėl $u(19-u)=84$, t. y. $u^2-19u+84=0$, $(u-7)(u-12)=0$. Kraštinės: $7\ \mathrm{cm}$ ir $12\ \mathrm{cm}$.
\kriterijai\begin{itemize}
\balas{1}{Gauna $u+v=19$ ir $uv=84$.}
\balas{1}{Sudaro kvadratinę lygtį arba pateikia pagrįstą sumos ir sandaugos paiešką.}
\balas{1}{Teisingai išsprendžia: $u=7$ arba $u=12$.}
\balas{1}{Pateikia abiejų kraštinių ilgius su vienetais arba su aiškia nuoroda, kad ilgiai centimetrais.}
\end{itemize}
\textit{Vertinimo pastaba.} Tinka ir pagrįstas skaičių poros parinkimas su abiejų sąlygų patikra; kvadratinės formulės naudoti nebūtina.
\newpage
\section*{Apolono variantas: 5--6 uždaviniai}
\uzd{5}{5}{$x^2-6x+4=0$; sudaryti redukuotą lygtį su šaknimis $x_1^2,x_2^2$, šaknų neskaičiuojant.}
\textbf{Sprendimas.} $S=6$, $P=4$. Nauja suma $S^2-2P=28$, nauja sandauga $P^2=16$. Atsakymas: $y^2-28y+16=0$.
\kriterijai\begin{itemize}
\balas{1}{Pagal Vijeto teoremą nustato $S=6$.}
\balas{1}{Nustato $P=4$.}
\balas{1}{Užrašo naujos sumos formulę $S^2-2P$.}
\balas{1}{Nustato naują sandaugą $P^2=16$.}
\balas{1}{Gauna naują sumą $28$ ir užrašo $y^2-28y+16=0$.}
\end{itemize}
\textit{Vertinimo pastaba.} Kintamąjį galima žymėti ir $x$. Vien apskaičiuojant pradines šaknis ir naudojant jas -- iki 2 t. (nauja sandauga ir galutinė lygtis).
\uzd{6}{5}{$(-x^2+6x-5)(x-2)(x-4)=-40$.}
\textbf{Sprendimas.} Imame $t=(x-3)^2$. Tada $-x^2+6x-5=4-t$ ir $(x-2)(x-4)=t-1$. Todėl $(4-t)(t-1)=-40$, t. y. $t^2-5t-36=0$, $(t-9)(t+4)=0$. Gauname $t=9$ arba $t=-4$. Kadangi $t=(x-3)^2\ge0$, reikšmę $-4$ atmetame. Iš $(x-3)^2=9$ gauname $x=0$ arba $x=6$.
\kriterijai\begin{itemize}
\balas{1}{Parenka keitinį ir gauna $(4-t)(t-1)=-40$ arba lygiavertę lygtį.}
\balas{1}{Gauna $t^2-5t-36=0$.}
\balas{1}{Randa abi reikšmes $t=9$ ir $t=-4$.}
\balas{1}{Pagrindžia $t\ge0$, atmeta $-4$ ir grįžta prie $(x-3)^2=9$.}
\balas{1}{Pateikia abu sprendinius $x=0$ ir $x=6$.}
\end{itemize}
\textit{Vertinimo pastaba.} Galima imti ir $u=x^2-6x$: tada $(u+5)(u+8)=40$, todėl $u(u+13)=0$. Reikšmė $u=0$ duoda $x=0,6$, o $u=-13$ realiųjų sprendinių neduoda, nes $(x-3)^2=-4$. Taškai skiriami už atitinkamus lygiaverčius žingsnius.

\newpage
\section*{Apolono variantas: 7--8 uždaviniai}
\uzd{7}{7}{$x^2-(m+2)x+2m=0$ turi skirtingas realiąsias šaknis, kurias galima pažymėti taip, kad $x_1^2=x_2$.}
\textbf{Sprendimas.} Imame $x_1=t$, $x_2=t^2$. Tada $t+t^2=m+2$, $t^3=2m$. Pašalinę $m$ gauname $t^3-2t^2-2t+4=0$, t. y. $(t-2)(t^2-2)=0$. Taigi $t=2,\sqrt2,-\sqrt2$. Iš $m=t^2+t-2$ gauname $m=4,\sqrt2,-\sqrt2$. Šaknų poros: $(2,4)$, $(\sqrt2,2)$, $(-\sqrt2,2)$; visos realiosios ir skirtingos.
\kriterijai\begin{itemize}
\balas{1}{Parenka $x_1=t$, $x_2=t^2$.}
\balas{1}{Užrašo sumą $t+t^2=m+2$.}
\balas{1}{Užrašo sandaugą $t^3=2m$.}
\balas{1}{Pašalina $m$ ir gauna kubinę lygtį.}
\balas{1}{Išskaido ir gauna visas tris $t$ reikšmes.}
\balas{1}{Gauna visas $m$ reikšmes: $4,\sqrt2,-\sqrt2$.}
\balas{1}{Patikrina arba aiškiai pagrindžia, kad visos poros yra skirtingos realiosios šaknys.}
\end{itemize}
\textit{Vertinimo pastaba.} Vien diskriminanto teigiamumas nepakeičia ryšio $x_1^2=x_2$ patikros. Pilnas ekvivalentus kelias ir nurodytos tinkamos poros pakanka.
\uzd{8}{7}{a) $mn-m-n+1$; b) rasti sveikuosius $a$, kad $x^2-ax+a+5=0$ turėtų dvi skirtingas sveikąsias šaknis.}
\textbf{Sprendimas.} \textbf{a)} $(m-1)(n-1)$.\par\textbf{b)} $x_1+x_2=a$, $x_1x_2=a+5$, todėl $(x_1-1)(x_2-1)=6$. Neatsižvelgiant į tvarką, daugiklių poros: $(1,6),(2,3),(-1,-6),(-2,-3)$. Šaknų poros: $(2,7),(3,4),(0,-5),(-1,-2)$. Atsakymas: $a\in\{-5,-3,7,9\}$.
\kriterijai\begin{itemize}
\balas{2}{a): 1 t. už grupavimą $m(n-1)-(n-1)$; 1 t. už $(m-1)(n-1)$. Teisinga pilna sandauga iš karto -- 2 t.}
\balas{1}{b): užrašo abi Vijeto lygybes.}
\balas{1}{b): gauna $(x_1-1)(x_2-1)=6$.}
\balas{1}{b): išvardija visas keturias daugiklių poras, įskaitant neigiamas.}
\balas{1}{b): iš jų gauna visas keturias $a$ reikšmes.}
\balas{1}{b): pagrindžia, kad šaknys sveikosios, skirtingos ir visos gautos reikšmės tinka.}
\end{itemize}
\textit{Vertinimo pastaba.} a) ir b) yra atskiri klaidų taisyklės vienetai; b) maksimumas -- 5 t. Patikrai pakanka nurodyti šaknų poras ir jų sumos bei sandaugos atitikimą.
\end{document}
